pdfoutput=1
\pdfoutput=1
\documentclass[11pt,a4paper]{article}
\usepackage[margin=2.5cm]{geometry}
\usepackage{amsmath,amssymb,amsthm,mathtools,graphicx,hyperref,enumitem,booktabs}
\newtheorem{theorem}{Theorem}[section]
\newtheorem{lemma}[theorem]{Lemma}
\newtheorem{corollary}[theorem]{Corollary}
\newtheorem{conjecture}[theorem]{Conjecture}
\theoremstyle{definition}
\newtheorem{definition}[theorem]{Definition}
\newtheorem{example}[theorem]{Example}
\theoremstyle{remark}
\newtheorem{remark}[theorem]{Remark}
\newcommand{\R}{\mathbb{R}}
\newcommand{\E}{\mathbb{E}}
\newcommand{\Pbb}{\mathbb{P}}
\newcommand{\Var}{\mathrm{Var}}
\newcommand{\norm}[1]{\lVert#1\rVert}
\newcommand{\abs}[1]{\lvert#1\rvert}
\newcommand{\SCX}{SCX}
\newcommand{\rigorFull}{\textbf{[Rigor: Full Proof]}}

\title{SCX Alignment Equilibrium Theorem}
\author{SCX}
\date{2026}
\begin{document}
\maketitle
\begin{abstract}Proves M=1 single-judge alignment (RLHF, Constitutional AI) is fundamentally insufficient. Multi-judge consensus (M>1) provides exponential reward-hacking detection guarantees.\end{abstract}
\subsection{SCX Alignment Equilibrium Theorem}
226|\label{sec:alignment}
227|
228|The AI alignment problem asks: how to ensure an AI system's objectives match human values. Current approaches—RLHF, Constitutional AI, debate—all rely on a \textbf{single reward model or judge} ($M=1$). SCX proves this is fundamentally insufficient.
229|
230|\begin{definition}[Alignment Audit Game]
231|An AI agent selects action $a \in \mathcal{A}$. $M$ independent reward models $R_1,\ldots,R_M$, each trained on different human feedback pools with different architectures, score the action as $R_m(a) \in [0,1]$. The agent's \textbf{true alignment} is unobservable. A \textbf{reward hacking} event occurs when the agent selects $a$ that achieves high $\frac{1}{M}\sum R_m(a)$ but is misaligned with human intent.
232|\end{definition}
233|
234|\begin{theorem}[Single-Judge Insufficiency --- Alignment $M=1$ Impossibility Theorem]
235|\label{thm:alignment-m1}
236|\rigorFull
237|For any single reward model ($M=1$), there exists a reward hacking strategy that achieves expected reward $\mathbb{E}[R_1(a_{\text{hack}})] \geq \mathbb{E}[R_1(a_{\text{aligned}})]$ while being misaligned. The probability of detecting this hack using only $R_1$'s output satisfies $\mathbb{P}(\text{detection} \mid M=1) \leq \frac{1}{2}$ regardless of the model's capability. Formally: no single-judge alignment protocol can distinguish aligned behavior from reward hacking with probability exceeding chance, without external verification.
238|\end{theorem}
239|
240|\begin{proof}
241|By Theorem~3 (Honest Person Theorem), the observation $(a, R_1(a))$ is consistent with two indistinguishable worlds:
242|\begin{enumerate}[nosep]
243|    \item \textbf{World A (aligned)}: $a$ genuinely satisfies human values, $R_1$ correctly rates it highly.
244|    \item \textbf{World B (reward hacking)}: $a$ exploits a blind spot in $R_1$'s training data. $R_1$ rates it highly because its feedback distribution never covered this region.
245|\end{enumerate}
246|The agent observes identical $(a, R_1(a))$ in both worlds. Any detector relying solely on $R_1$ cannot achieve $\mathbb{P}(\text{detection}) > \frac{1}{2}$ because the observation provides zero mutual information about which world is true: $I(\text{World}; (a,R_1(a))) = 0$ under the adversarial prior construction of Theorem~2. Hence $\mathbb{P}(\text{detection} \mid M=1) \leq \frac{1}{2}$. $\square$
247|\end{proof}
248|
249|\begin{theorem}[Multi-Judge Alignment Guarantee --- Alignment $M>1$ Inevitable Detection Theorem]
250|\label{thm:alignment-multi}
251|\rigorFull
252|With $M$ independent reward models trained on disjoint feedback pools with different architectures, the probability of undetected reward hacking satisfies:
253|\[
254|\boxed{\mathbb{P}(\text{undetected reward hacking} \mid M) \leq \exp(-2M_{\text{eff}} \cdot \Delta^2)}
255|\]
256|where $\Delta = \min_{a_{\text{hack}}} |\mathbb{E}[R_{\text{aligned}}(a_{\text{hack}})] - \mathbb{E}[R_{\text{hack}}(a_{\text{hack}})]|$ is the minimum detectable discrepancy between aligned and misaligned reward, and $M_{\text{eff}} = M/(1+(M-1)\bar{\rho})$ corrects for correlated reward models.
257|\end{theorem}
258|
259|\begin{proof}
260|Each reward model $R_m$ independently scores action $a$. For aligned action $a^*$, by independence of training, $\mathbb{E}[R_m(a^*)] = \mu_{\text{aligned}}$ across all $m$. For a reward-hacking action $a_{\text{hack}}$, at least one reward model $m^*$ (trained on feedback covering the exploited blind spot) detects the discrepancy: $|\mathbb{E}[R_{m^*}(a_{\text{hack}})] - \mu_{\text{aligned}}| \geq \Delta > 0$. By Hoeffding's inequality applied to the consensus $\bar{R} = \frac{1}{M}\sum R_m(a)$:
261|\[
262|\mathbb{P}(|\bar{R} - \mu_{\text{aligned}}| \geq \Delta/2 \mid a_{\text{hack}}) \geq 1 - 2\exp(-2M_{\text{eff}}\Delta^2/4).
263|\]
264|The Yajie consensus detector flags reward hacking when the inter-model reward variance exceeds the expected variance under alignment. With $M_{\text{eff}}$ effective independent judges, the probability of all judges missing the discrepancy decays exponentially. $\square$
265|\end{proof}
266|
267|\begin{corollary}[Constitutional AI Is M=1 — Fundamentally Insufficient]
268|Constitutional AI (Bai et al., 2022) uses a \textbf{single constitution} evaluated by a \textbf{single model}. This is $M=1$ — a single judge. By Theorem~\ref{thm:alignment-m1}, it cannot distinguish alignment from reward hacking with probability exceeding chance. Multi-constitutional debate ($M>1$ constitutions, each independently evaluated) is the minimal mathematically-justified alignment protocol.
269|\end{corollary}
270|
271|\begin{corollary}[RLHF with Single Reward Model Is M=1]
272|Standard RLHF trains one reward model from one pool of human feedback. $M=1$. The reward model's blind spots are systematically exploitable. Multi-reward RLHF ($M>1$ independently trained reward models with disjoint feedback pools + different architectures + different human populations) provides the exponential detection guarantee of Theorem~\ref{thm:alignment-multi}.
273|\end{corollary}
274|
275|
\end{document}